% Estilo común de las presentaciones. Los datos (título, autor…) salen
% de ../memoria/config/datos.tex, igual que en la memoria.
% =====================================================================
%  DATOS DEL TRABAJO
%  Es el ÚNICO fichero que hay que editar al empezar. Todo lo demás
%  (portada, metadatos del PDF, presentaciones, nombre de los
%  entregables generados por la CI) se alimenta de estas macros.
% =====================================================================

% --- Tipo de trabajo -------------------------------------------------
\newcommand{\TipoTrabajo}{Trabajo Fin de Grado}      % o: Trabajo Fin de Máster
\newcommand{\TipoCorto}{TFG}                         % o: TFM
\newcommand{\Titulacion}{Grado en Ingeniería en Tecnologías Industriales}

% --- Título (español e inglés) ---------------------------------------
\newcommand{\Titulo}{Título del trabajo en español}
\newcommand{\TituloEN}{Project title in English}

% --- Personas ---------------------------------------------------------
\newcommand{\Autor}{Nombre Apellido1 Apellido2}
\newcommand{\Director}{Nombre Apellido1 Apellido2}
\newcommand{\Codirector}{}                           % dejar vacío si no hay

% --- Institución ------------------------------------------------------
\newcommand{\Universidad}{Universidad Pontificia Comillas}
\newcommand{\Escuela}{Escuela Técnica Superior de Ingeniería (ICAI)}
\newcommand{\LogoUniversidad}{figuras/logo_universidad}

% --- Fechas -----------------------------------------------------------
\newcommand{\Lugar}{Madrid}
\newcommand{\FechaEntrega}{Junio 2027}
\newcommand{\CursoAcademico}{2026-2027}

% --- Palabras clave -----------------------------------------------------
\newcommand{\PalabrasClave}{palabra 1, palabra 2, palabra 3, palabra 4}
\newcommand{\Keywords}{keyword 1, keyword 2, keyword 3, keyword 4}

% --- Entrega ----------------------------------------------------------
% Nombre base de los PDF que publica la CI (sin espacios ni tildes).
% Convención: <TIPO>_<Apellido1><Apellido2>_<Nombre>
\newcommand{\NombreEntrega}{TFG_Apellido1Apellido2_Nombre}

% Anexo I (declaración de originalidad, uso de IA y autorización del
% director). La normativa exige: portada -> Anexo I firmado -> portada.
% Poner \AnexoItrue cuando el PDF firmado esté en la ruta indicada.
\newif\ifAnexoI
\AnexoIfalse
\newcommand{\RutaAnexoI}{administrativo/AnexoI_firmado.pdf}


\usepackage[utf8]{inputenc}
\usepackage[T1]{fontenc}
\usepackage[spanish, es-tabla, es-noshorthands]{babel}
\usepackage{etoolbox}
\usepackage{booktabs}
\usepackage{siunitx}
\sisetup{output-decimal-marker={,}}
\usepackage{tikz}
\usetikzlibrary{arrows.meta, positioning}

% Reutiliza las figuras de la memoria sin duplicarlas
\graphicspath{{../memoria/}{../memoria/figuras/}}

% metropolis usa Fira con XeLaTeX/LuaLaTeX; con pdflatex cae a la fuente
% por defecto sin problema, así que se silencia ese aviso.
\usepackage{silence}
\WarningFilter{beamerthememetropolis}{You need to compile with XeLaTeX}
\usetheme{metropolis}
\metroset{progressbar=frametitle, numbering=fraction, block=fill}
\definecolor{acento}{HTML}{1E3A5F}
\setbeamercolor{frametitle}{bg=acento}
\setbeamercolor{progress bar}{fg=acento!60}
\setbeamercolor{alerted text}{fg=acento!80!red}

\title{\Titulo}
\author{\texorpdfstring{\Autor\\{\small Director: \Director\ifdefempty{\Codirector}{}{ · Codirector: \Codirector}}}{\Autor}}
\institute{\footnotesize \TipoTrabajo · \Titulacion\\\Escuela · \Universidad}
% Logo en la esquina de la portada (con tikz para no descuadrar la página)
\usetikzlibrary{calc}
\addtobeamertemplate{title page}{}{%
  \begin{tikzpicture}[remember picture, overlay]
    \node[anchor=north east] at ($(current page.north east)+(-0.5,-0.4)$)
      {\includegraphics[height=1.1cm]{\LogoUniversidad}};
  \end{tikzpicture}}
